\documentclass[11pt,a4paper, 
swedish, english %% Make sure to put the main language last!
]{article}
\pdfoutput=1

%% Andréas's custom package 
%% (Will work for most purposes, but is mainly focused on physics.)
\usepackage{../custom_as}

%% Figures can now be put in a folder: 
\graphicspath{ {figures/} %{some_folder_name/}
}

%% If you want to change the margins for just the captions
\usepackage[margin=10 pt]{caption}

%% To add todo-notes in the pdf
\usepackage[%disable  %%this will hide all notes
]{todonotes} 

%% Change the margin in the documents
\usepackage[
%            top    = 3cm,              %% top margin
%            bottom = 3cm,              %% bottom margin
%            left   = 2.5cm, right  = 2.5cm %% left and right margins
]{geometry}


%% If you want to change the formatting of the section headers
%\renewcommand{\thesection}{...}
\renewcommand{\thesubsection}{\arabic{section} (\alph{subsection})}


%%%%%%%%%%%%%%%%%%%%%%%%%%%%%%%%%%%%%%%%%%%%%%%%%%%%%%%%%%%%%%%%%%%%%%
\begin{document}%% v v v v v v v v v v v v v v v v v v v v v v v v v v
%%%%%%%%%%%%%%%%%%%%%%%%%%%%%%%%%%%%%%%%%%%%%%%%%%%%%%%%%%%%%%%%%%%%%%


%% If you want to use an external file for the title page
%\renewcommand{\thefootnote}{\fnsymbol{footnote}}

%kortkommandon för mailaddresserna
\newcommand{\andsunds}{andsunds@student.chalmers.se}
\newcommand{\rigon}{rigon@student.chalmers.se}



\pagenumbering{roman} %%Romersk sidnumrering i början
\begin{titlepage}
\newgeometry{top=3cm, bottom=2cm}

\newcommand{\HRule}{\rule{\linewidth}{0.5mm}} % Defines a new command for the horizontal lines, change thickness here

\center % Center everything on the page
 
%------------------------------------------------------------------------------------
%	HEADING SECTIONS
%------------------------------------------------------------------------------------

\textsc{\huge Chalmers tekniska högskola}\\[1.5cm] % Name of university/college
\textsc{\Large Rapport, Experimentell fysik 2}\\[0.2cm] % Major heading such as course name
\textsc{\large Termodynamik -- Uppgift 3 }\\[0.5cm] % Minor heading such as course title

%------------------------------------------------------------------------------------
%	TITLE SECTION
%------------------------------------------------------------------------------------

\HRule \\[0.4cm]
{ \LARGE \bfseries 
Studier av kvicksilveratomens atomära emissionsspektra samt absorptionsspektra av laserfärgämnena Rhodamin B och Kumarin 307
}\\[0.4cm] % Title of  document
\HRule \\[1.5cm]
 
%------------------------------------------------------------------------------------
%	AUTHOR SECTION
%------------------------------------------------------------------------------------

\begin{minipage}{0.4\textwidth}
\begin{flushleft} \large
\emph{Författare:}\\
Andréas Sundström\footnotemark{} \\
Rigon Demisai\footnotemark{} 
\end{flushleft}
\end{minipage}
~
\begin{minipage}{0.4\textwidth}
\begin{flushright} \large
\emph{Labassistent:} \\
Martin Wersäll
\end{flushright}
\end{minipage}\\[3cm]

\setcounter{footnote}{0}
\stepcounter{footnote}
  \footnotetext{\href{mailto:\andsunds}{\texttt{\andsunds}}}
\stepcounter{footnote}
  \footnotetext{\href{mailto:\rigon}{\texttt{\rigon}}}



%------------------------------------------------------------------------------------
%	DATE SECTION
%------------------------------------------------------------------------------------
% Följer ISO-standarden för tidsintervall:
% https://en.wikipedia.org/wiki/ISO_8601#Time_intervals
% "Double hyphen" också ok istället för '/'. -- i LaTeX är dock lite på gränsen
{ \large
\begin{tabular}{rc}
    Laboration utförd: & 2015-12-11/15 \\[0.1cm]
    Rapport inlämnad: & \today
\end{tabular}\\[1cm]
}

%------------------------------------------------------------------------------------
%	LOGO SECTION
%------------------------------------------------------------------------------------

\includegraphics[height=5cm]{logo.pdf} % Include a department/university logo
 
%------------------------------------------------------------------------------------

\vfill % Fill the rest of the page with whitespace

\end{titlepage}
\restoregeometry


\setcounter{page}{2}%detta är ANDRA (2) sidan

\renewcommand{\abstractname}{Sammandrag}
\begin{abstract}
Den här rapporten beskriver en spektroskopisk studie av
kvicksilveratomens atomära emissionsspektrum samt en studie av
laserfärgämnena rhodamin~B och kumarin~307 och deras
absorptionsspektra.
Ur emisionspektrumet från Hg har vissa av atomens energinivåer kartlagts,
baserat på kvanmekaniska regler och jämförelse med tidigare data på
Hg.  Kvicksilvrets emissionsspektra
är taget i intervallet 350--1100\,nm, där detekterades totalt 23
emissionstoppar, varefter 13 övergångar kunde identifieras.
Absorptionsspektrumen för rhodamin~B och kumarin~307 visar breda
absorptionsband vilket är kännetecknande för flourescerande ämnen som
består av stora organiska molekyler. 
Mätningarna har utförts med en Spex~270M spektrometer och
datainsamlingen har gjorts i LabView.
\end{abstract}

\renewcommand{\abstractname}{Abstract}
\begin{abstract}
This report describes a spectroscopic study of the atomic emision
spectrum of mercury, and also a study of the laser dyes Rhodamine~D
and Coumarin~307 and their absorption spectrum.
Som of the energylevels of Hg have been idetified from the emission
spectrum, based on quantum mecanichal rules and comparison with erlier
data on Hg. The emission spectrum of mercury was taken in the interval
of 350--1100\,nm, there a total of 23 emission lines were detected,
from which 13 different atomic transitions could be identified. 
The absorption spectrum of Rhodamin~B and Coumarin~307  show wide
absorption bands which are characteristic of big organic molecules.
THe measurements were made with a Spex~270M spectrometer and the data
collection was done through LabVIEW.
\end{abstract}

\clearpage
\renewcommand{\contentsname}{Innehållsförteckning}
\tableofcontents

\clearpage
\pagenumbering{arabic}
\setcounter{page}{1}

\renewcommand{\thefootnote}{\arabic{footnote}}
\setcounter{footnote}{0}



%%%%%%%%%%%%%%%%%%%% vvv Internal title page vvv %%%%%%%%%%%%%%%%%%%%%
\title{Gravitation and Cosmology -- FFM 071
\\ {\Large Hand-in set 1} }
\author{Andréas Sundström}
\date{\today}

\maketitle

%%%%%%%%%%%%%%%%%%%% ^^^ Internal title page ^^^ %%%%%%%%%%%%%%%%%%%%%
%% If you want a list of all todos
%\todolist

\newcommand{\EM}{\text{EM}}

\section{Some Lorentz transformations}


\subsection{Proper time}
To show that the proper time is invariant, we just use the definition
of $\rd\tau$ and the transformed coordinates 
$\rd{x'}^\alpha={\Lambda^\alpha}_\gamma \rd{x}^{\gamma}$,
\begin{equation}
\rd{\tau'}^2=-\eta_{\alpha\beta}\rd{x'}^{\alpha}\rd{x'}^{\beta}
=-\eta_{\alpha\beta}({\Lambda^\alpha}_\gamma \rd{x}^{\gamma})
({\Lambda^\beta}_\delta\rd{x}^{\delta})
=-\eta_{\alpha\beta}{\Lambda^\alpha}_\gamma{\Lambda^\beta}_\delta
\rd{x}^{\gamma}\rd{x}^{\delta}.
\end{equation}
next we use the condition that the Lorentz transform must satisfy
\begin{equation}
\eta_{\alpha\beta}{\Lambda^\alpha}_\gamma{\Lambda^\beta}_\delta
=\eta_{\gamma\delta},
\end{equation}
which therefore yields
\begin{equation}
\rd{\tau'}^2=\eta_{\gamma\delta}\rd{x}^{\gamma}\rd{x}^{\delta}
=\rd\tau^2.
\end{equation}
Since we have now concluded that $\rd\tau^2$ is a scalar (invariant)
under Lorentz transforms, so must $\rd\tau=\sqrt{\rd\tau^2}$.

\subsection{Relativistic particle velocity}
For a Lorentz boost with $\vb*v=(v,0,0)$, the transformations are as
follows
\begin{equation}
\begin{cases}
\rd{t}'=\gamma(\rd{t}-v\rd{x})\\
\rd{x}'=\gamma(\rd{x}-v\rd{t})\\
\rd{y}'=\rd{y}\\
\rd{z}'=\rd{z}
\end{cases}
\qand
\begin{cases}
\rd{t}=\gamma(\rd{t}'+v\rd{x}')\\
\rd{x}=\gamma(\rd{x}'+v\rd{t}')\\
\rd{y}=\rd{y}'\\
\rd{z}=\rd{z}'.
\end{cases}
\end{equation}
The velocity of a particle 
\begin{equation}
\vb*u=\qty(\dv{x}{t},\dv{y}{t},\dv{z}{t})
=(u_x,u_y,u_z)=(u\cos\theta,u\sin\theta,0)
\end{equation}
in the first frame, will transform as
\begin{equation}\label{eq1:upx1}
u'_x=\dv{x'}{t'}=\dv{x'}{t}\dv{t}{t'}
=\qty(\pdv{x'}{t}+\pdv{x'}{x}\pdv{x}{t})
\qty(\pdv{t}{t'}+\pdv{t}{x'}\pdv{x'}{t'})
=\gamma^2(-v+u_x)(1+vu'_x)
\end{equation}
and
\begin{equation}\label{eq1:upy1}
u'_y=\dv{y'}{t'}=\dv{y'}{t}\dv{t}{t'}
=\qty(\pdv{y'}{y}\pdv{y}{t})
\qty(\pdv{t}{t'}+\pdv{t}{x'}\pdv{x'}{t'})
=\gamma u_y(1+vu'_x).
\end{equation}
Solving for $u'_x$ in \eqref{eq1:upx1} yields
\begin{equation}
u'_x=\frac{u_x-v}{\gamma^{-2}-v(u_x-v)}
=\frac{u_x-v}{1-v^2-v(u_x-v)}
=\frac{u_x-v}{1-u_x v},
\end{equation}
similarly with \eqref{eq1:upy1}
\begin{equation}
u'_y=\gamma u_y\qty(1+v\frac{u_x-v}{1-u_x v})
=\frac{u_y}{\gamma(1-u_x v)}.
\end{equation}
The transform of $\theta$, the angle $\vb*u$ makes with the $x$ axis,
is easily given by 
\begin{equation}
\tan\theta'=\frac{u_y'}{u_x'}
=\qty(\frac{u\sin\theta}{\gamma(1-u\cos\theta v)})
\qty(\frac{1-u\cos\theta v}{u\cos\theta-v})
=\frac{\sin\theta}{\gamma(\cos\theta-\nicefrac{v}{u})}.
\end{equation}



\section{The stress tensor}
To show the conservation laws for the stress tensor\footnotemark{}
\footnotetext{I call the stress tensor of a particle
  ${T_n}^{\alpha\beta}$, with subscript $n$, and use $x_n(\tau)$ and
  $p_n$ for the trajectory and momentum of the particle, to more
  clearly distinguish what ``belongs'' to the particle and what are
  free variables.} 
\begin{equation}
{T_n}^{\alpha\beta}(x)=\int\rd\tau\,{p_n}^\alpha
\dv{{x_n}^\beta}{\tau}\delta^4\big(x-x_n(\tau)\big).
\end{equation}
The divergence of the stress tensor therefore becomes
\begin{equation}\label{eq2:pdvT1}
\pd_\beta{T_n}^{\alpha\beta}(x)=\int\rd\tau\,{p_n}^\alpha
\dv{{x_n}^\beta}{\tau}\,\pdv{x^\beta}\delta^4\big(x-x_n(\tau)\big).
\end{equation}
We note that the derivatives only act on the delta function since
$x_n$ and $p_n$ are the trajectory and momentum of the particle, while
the derivatives act on the free variable $x$. Next we note that for
any function $f$ 
\begin{equation}
\pdv{x}\Big[f(x-y)\Big]=-\pdv{y}\Big[f(x-y)\Big],
\end{equation}
which means that \eqref{eq2:pdvT1} can be written as
\vspace{-1em}
\begin{equation}
\begin{aligned}
\pd_\beta{T_n}^{\alpha\beta}(x)=&-\int\rd\tau\,{p_n}^\alpha
\overbrace{\dv{{x_n}^\beta}{\tau}\,\pdv{{x_n}^\beta}}^{=\dv*\tau}
\delta^4\big(x-x_n(\tau)\big)\\
=&-\int\rd\tau\,{p_n}^\alpha\dv{\tau}\delta^4\big(x-x_n(\tau)\big).
\end{aligned}
\end{equation}
Integration by parts now yields
\begin{equation}
\begin{aligned}
\pd_\beta{T_n}^{\alpha\beta}(x)=&
-\int\rd\tau\,{p_n}^\alpha\dv{\tau}\delta^4\big(x-x_n(\tau)\big)\\
=&-\cancel{\qty[{p_n}^\alpha\delta^4\big(x-x_n(\tau)\big)]}
+\int\rd\tau\,\dv{{p_n}^\alpha}{\tau}\delta^4\big(x-x_n(\tau)\big).
\end{aligned}
\end{equation}
The first term from the integration by parts cancel since the
integration is over ``all $\tau$'' and therefore $x_n\neq x$ at the
end points, which means that the delta function is zero.

Now if the particle interacts with an electromagnetic field, it will
experience a Lorentz force
\begin{equation}
\dv{{p_n}^\alpha}{\tau}=e_n\eta_{\beta\gamma}F^{\alpha\beta}
\dv{{x_n}^\gamma}{\tau}={F^{\alpha}}_\gamma \,e_n{u_n}^\gamma.
%={F^{\alpha}}_\gamma J^\gamma.
\end{equation}
This means that
\begin{equation}\label{eq2:pdvT}
\pd_\beta{T_n}^{\alpha\beta}(x)=\int\rd\tau\,{F^{\alpha}}_\gamma e_n{u_n}^\gamma
\delta^4\big(x-x_n(\tau)\big)
={F^{\alpha}}_\gamma e_n{u_n}^\gamma
\delta^3\big(\vb*x-\vb*x_n(\tau)\big)
={F^{\alpha}}_\gamma J^\gamma
\end{equation}
which in general is non-vanishing. Note that since we are assuming
that this particle is the only particle the current is indeed
$J^\alpha=e_n{u_n}^\alpha\delta^3\big(\vb*x-\vb*x_n(\tau)\big)$, and
if we were to have more particles all the above calculations would be
exactly the same (but with summation signs in front of them). 

We now take a look at the electromagnetic stress tensor
\begin{equation}
{T_{\EM}}^{\alpha\beta}={F^{\alpha}}_\gamma F^{\beta\gamma}
-\frac{1}{4}\eta^{\alpha\beta}F_{\gamma\delta}F^{\gamma\delta},
\end{equation}
and then Maxwell's equations will come in handy later
\begin{subequations}
\begin{equation}\label{eq:Maxwell}
\pd_\alpha F^{\alpha\beta}=-J^\beta,
\end{equation}
\begin{equation}\label{eq:Bianchi}
%\pd_{[\alpha}F_{\beta\gamma]}=
\pd_\alpha F_{\beta\gamma} +\pd_\beta F_{\gamma\alpha} +\pd_\gamma F_{\beta\alpha} 
=0.
\end{equation}
\end{subequations}
Its divergence therefore is
\begin{equation}
\begin{aligned}
\pd_{\beta}{T_{\EM}}^{\alpha\beta}=&
{F^{\alpha}}_\gamma \qty(\pd_\beta F^{\beta\gamma})
+\qty(\pd_\beta{F^{\alpha}}_\gamma) F^{\beta\gamma}
-\frac{1}{4}\eta^{\alpha\beta}\pd_\beta\qty[F_{\gamma\delta}F^{\gamma\delta}]\\
=&-{F^{\alpha}}_\gamma J^\gamma
+F^{\beta\gamma} \,\pd_\beta{F^{\alpha}}_\gamma
-\frac{1}{2}F_{\gamma\delta}\,\pd^\alpha F^{\gamma\delta}.
\end{aligned}
\end{equation}
The first term is given directly from \eqref{eq:Maxwell}.
We now take a look at the last two terms. Firstly we can raise and
lower the betas and gammas which are contracted in the second term, 
$F^{\beta\gamma} \,\pd_\beta{F^{\alpha}}_\gamma
=F_{\beta\gamma} \,\pd^\beta F^{\alpha\gamma}$, then in the last term
we are free to rename the two contracted indices $\gamma\to\beta$ and
$\delta\to\gamma$. We now have
\begin{equation}
F_{\beta\gamma} \,\pd^\beta F^{\alpha\gamma}
-\frac{1}{2}F_{\beta\gamma}\,\pd^\alpha F^{\beta\gamma}=
\frac{1}{2}\qty( F_{\beta\gamma}\pd^\beta F^{\alpha\gamma}
-F_{\beta\gamma}\pd^\beta F^{\gamma\alpha})
-\frac{1}{2}F_{\beta\gamma}\pd^\alpha F^{\beta\gamma}.
\end{equation}
Here we used the anti-symmetry of the EM tensor, 
$F^{\alpha\gamma}=-F^{\gamma\alpha}$. Now we do some more index
gymnastics by interchanging $\beta\leftrightarrow\gamma$ in the first
term, $F_{\gamma\beta}\pd^\gamma F^{\alpha\beta}
=-F_{\beta\gamma}\pd^\gamma F^{\alpha\beta}$. 
In the end what we have is
\begin{equation}
F_{\beta\gamma} \,\pd^\beta F^{\alpha\gamma}
-\frac{1}{2}F_{\beta\gamma}\,\pd^\alpha F^{\beta\gamma}=
-\frac{1}{2}F_{\beta\gamma}\qty( 
\pd^\gamma F^{\alpha\beta}
+\pd^\beta F^{\gamma\alpha}
+\pd^\alpha F^{\beta\gamma}).
\end{equation}
The expression inside the parentheses is just the contravariant version
of the RHS of \eqref{eq:Bianchi} and therefore vanishes.
So what we are left with for the divergence is
\begin{equation}
\pd_{\beta}{T_{\EM}}^{\alpha\beta}=-{F^{\alpha}}_\gamma J^\gamma,
\end{equation}
which is exactly $-\pd_\beta{T_n}^{\alpha\beta}$ from \eqref{eq2:pdvT}.

Therefore since $\pd_\beta$ is a linear operator we know that
\begin{equation}
\pd_{\beta}\qty[{T_{\EM}}^{\alpha\beta}+{T_n}^{\alpha\beta}]=0.
\end{equation}
This is physically reasonable since the total stress tensor, which
represents the total energy and momentum-current density, should
reasonably be conserved in a system. The two terms
${T_{\EM}}^{\alpha\beta}$ and ${T_n}^{\alpha\beta}$ are not conserved
separately since the particle and EM field exchanges energy and
momentum with each other, which also leads to the above conclusion
that there is no \emph{net} change in the total energy or momentum of
the system.







%%%%%%%%%%%%%%%%%%%%%%%%%% The bibliography %%%%%%%%%%%%%%%%%%%%%%%%%%
%\newpage
%% This bibliography uses BibTeX
\bibliographystyle{ieeetr}
\bibliography{references}%requires a file named 'references.bib'
%% Citations are as usual: \cite{example_article}

%%%%%%%%%%%%%%%%%%%%%%%%%%%%% Appendices %%%%%%%%%%%%%%%%%%%%%%%%%%%%%
\clearpage %% on a new page 
\appendix  %% This will change the page numbering to A1, A2, A3, ...;
           %% and also change the sections to A, A.1, ...; B, B.1, ...


%%%%%%%%%%%%%%%%%%%%%%%%%%%%%%%%%%%%%%%%%%%%%%%%%%%%%%%%%%%%%%%%%%%%%%
\end{document}%% ^ ^ ^ ^ ^ ^ ^ ^ ^ ^ ^ ^ ^ ^ ^ ^ ^ ^ ^ ^ ^ ^ ^ ^ ^ ^ ^
%%%%%%%%%%%%%%%%%%%%%%%%%%%%%%%%%%%%%%%%%%%%%%%%%%%%%%%%%%%%%%%%%%%%%%




%%%%  Some (useful) templates


%% På svenska ska citattecknet vara samma i både början och slut.
%% Använd två apostrofer: ''.


%% Including PDF-documents
\includepdf[pages={1-}]{filnamn.pdf} % NO blank spaces in the file name

%% Figures (pdf, png, jpg, ...)
\begin{figure}\centering
\centerline{ % centers figures larges than 1\textwidth
\includegraphics[width=.8\textwidth]{file_name.pdf}
}
\caption{}
\label{fig:}
\end{figure}

%% Figures from xfig's "Combined PDF/LaTeX"
\begin{figure}\centering
\resizebox{.8\textwidth}{!}{\input{file_name.pdf_t}}
\caption{}
\label{fig:}
\end{figure}


%% If you want to add something to the ToC
%% (Without having an actual header in the text.)
\stepcounter{section} %For example a 'section'
\addcontentsline{toc}{section}{\Alph{section}\hspace{8 pt}Labblogg} 


%  LocalWords:  contravariant
